%==============================================================================
% Chapitre 28 - Moments d'inertie
%==============================================================================

\section{Introduction}

Dans les chapitres précédents, nous avons étudié :

\begin{itemize}
    \item la masse ;
    \item le centre de gravité ;
    \item le centre de pression.
\end{itemize}

Cependant, deux projectiles possédant la même masse totale et le même centre
de gravité peuvent réagir différemment à une perturbation.

Cette différence provient notamment de la manière dont la masse est répartie
dans le volume du projectile.

Cette propriété est décrite par les moments d'inertie.

\begin{aretenir}

Le moment d'inertie caractérise la résistance d'un corps à une modification
de son mouvement de rotation.

\end{aretenir}

\section{Analogie avec la masse}

En mouvement de translation :

\begin{itemize}
    \item la masse mesure la résistance à l'accélération.
\end{itemize}

En mouvement de rotation :

\begin{itemize}
    \item le moment d'inertie mesure la résistance à l'accélération
          angulaire.
\end{itemize}

Les deux notions jouent des rôles analogues.

\section{Exemple intuitif}

Considérons une haltère.

Si la masse est concentrée près du centre :

\begin{itemize}
    \item elle tourne facilement.
\end{itemize}

Si la masse est placée loin du centre :

\begin{itemize}
    \item il devient plus difficile de modifier sa rotation.
\end{itemize}

Pourtant, la masse totale peut être identique.

\section{Définition}

Pour un ensemble de masses discrètes :

\begin{equation}
I = \sum m_i r_i^2
\end{equation}

où :

\begin{itemize}
    \item $I$ représente le moment d'inertie ;
    \item $m_i$ représente une masse ;
    \item $r_i$ représente sa distance à l'axe de rotation.
\end{itemize}

Cette équation montre qu'une masse éloignée de l'axe contribue fortement au
moment d'inertie.

\section{Unités}

Dans le Système international, le moment d'inertie s'exprime en :

\begin{equation}
kg \cdot m^2
\end{equation}

Cette unité traduit la combinaison :

\begin{itemize}
    \item d'une masse ;
    \item d'une distance au carré.
\end{itemize}

\section{Axes de rotation}

Un projectile peut tourner autour de plusieurs axes.

On distingue notamment :

\begin{itemize}
    \item l'axe longitudinal ;
    \item les axes transverses.
\end{itemize}

Les moments d'inertie associés sont généralement différents.

\section{Moment d'inertie longitudinal}

Le moment d'inertie longitudinal correspond à une rotation autour de l'axe du
projectile.

C'est celui qui intervient dans la rotation imposée par les rayures.

Il influence notamment :

\begin{itemize}
    \item la vitesse de rotation ;
    \item la stabilité gyroscopique ;
    \item la conservation du spin.
\end{itemize}

\section{Moment d'inertie transverse}

Le moment d'inertie transverse correspond à une rotation autour d'un axe
perpendiculaire à l'axe du projectile.

Il intervient notamment lors :

\begin{itemize}
    \item de la précession ;
    \item de la nutation ;
    \item des changements d'orientation.
\end{itemize}

\section{Influence de la géométrie}

La forme du projectile influence fortement les moments d'inertie.

Par exemple :

\begin{itemize}
    \item un projectile long possède généralement un moment d'inertie
          transverse élevé ;
    \item un projectile court possède généralement un moment d'inertie
          transverse plus faible.
\end{itemize}

\section{Influence de la répartition des masses}

La répartition interne est également importante.

Deux projectiles de même géométrie extérieure peuvent présenter :

\begin{itemize}
    \item des moments d'inertie différents ;
    \item des comportements dynamiques différents.
\end{itemize}

\section{Moments principaux d'inertie}

Pour un projectile présentant une symétrie de révolution, on distingue
souvent :

\begin{itemize}
    \item un moment longitudinal ;
    \item deux moments transverses identiques.
\end{itemize}

Cette simplification est largement utilisée en balistique.

\section{Lien avec les moments aérodynamiques}

Les forces aérodynamiques génèrent des moments autour du centre de gravité.

La réaction du projectile dépend alors :

\begin{itemize}
    \item de l'intensité de ces moments ;
    \item des moments d'inertie du projectile.
\end{itemize}

Un même effort aérodynamique peut produire des effets très différents selon
le projectile considéré.

\section{Importance pour la stabilité}

Les moments d'inertie interviennent directement dans les équations de
stabilité.

Ils influencent notamment :

\begin{itemize}
    \item la précession ;
    \item la nutation ;
    \item les oscillations du projectile.
\end{itemize}

Ils constituent donc un paramètre fondamental des études balistiques
avancées.

\section{Moments d'inertie et gyroscope}

Un projectile en rotation se comporte partiellement comme un gyroscope.

Sa réaction aux perturbations dépend :

\begin{itemize}
    \item de sa vitesse de rotation ;
    \item de ses moments d'inertie.
\end{itemize}

Cette propriété sera étudiée dans les prochains chapitres.

\section{Mesure et calcul}

Les moments d'inertie peuvent être :

\begin{itemize}
    \item calculés à partir de la géométrie ;
    \item obtenus par modélisation numérique ;
    \item mesurés expérimentalement.
\end{itemize}

Les modèles modernes utilisent fréquemment des outils de CAO pour les
déterminer.

\section{Moments d'inertie et simulation}

Les modèles de masse ponctuelle ignorent les moments d'inertie.

Les modèles avancés les utilisent explicitement.

Ils deviennent indispensables dans les modèles :

\begin{itemize}
    \item 5DOF ;
    \item 6DOF ;
    \item dynamiques complets.
\end{itemize}

\begin{simulation}

La connaissance des moments d'inertie est nécessaire pour reproduire les
phénomènes de précession et de nutation dans une simulation.

\end{simulation}

\section{Vers la stabilité statique}

Nous disposons maintenant des éléments suivants :

\begin{itemize}
    \item géométrie ;
    \item masse ;
    \item centre de gravité ;
    \item centre de pression ;
    \item moments d'inertie.
\end{itemize}

Nous pouvons désormais répondre à une question essentielle :

\begin{quote}
Pourquoi un projectile revient-il naturellement vers sa position d'équilibre
ou, au contraire, devient-il instable ?
\end{quote}

Cette question conduit à la notion de stabilité statique.

\section{À retenir}

\begin{aretenir}

\begin{itemize}
    \item Le moment d'inertie est l'équivalent rotationnel de la masse.
    \item Il dépend de la répartition des masses.
    \item Un projectile possède plusieurs moments d'inertie.
    \item Ces grandeurs influencent la stabilité et les oscillations.
    \item Elles sont indispensables aux modèles balistiques avancés.
\end{itemize}

\end{aretenir}
